\documentclass[10pt,a4paper]{amsart}
\usepackage[margin=19mm]{geometry}
\usepackage{amsmath,amssymb,mathtools}
\usepackage[scr=rsfs,cal=euler]{mathalfa}
\usepackage{xfrac}
\usepackage[unicode,hidelinks,bookmarks=false]{hyperref}
\newcommand{\ie}{i.e.\ }
\newcommand{\cf}{cf.\ }
\newcommand{\resp}{respectively\ }
\newcommand{\Subsection}{\subsection}
\providecommand{\nobreakdash}{\nobreak\hbox{-}}
\setlength{\emergencystretch}{2em}
\multlinegap0pt
\usepackage{bm}
\usepackage{bbm}
\usepackage{dsfont}
\usepackage{xspace}
\usepackage{cancel}
\usepackage{mathtools}
\usepackage{amsthm}
\makeatletter
\@ifundefined{theorem}{\newtheorem{theorem}{Theorem}}{}
\@ifundefined{lemma}{\newtheorem{lemma}[theorem]{Lemma}}{}
\@ifundefined{proposition}{\newtheorem{proposition}[theorem]{Proposition}}{}
\makeatother
\def\G{\mathbb{G}}
\newcommand{\cP}{ \mathcal{P}}
\title[Master equations with an individual noise on finite state gr]{Master equations with an individual noise on finite state graphs}
\author{Wilfrid Gangbo, Sebastian Munoz, Jeremy Wu, Zhaoyu Zhang}
\date{Source: arXiv:2605.06589v1 (CC BY 4.0)}
\begin{document}
\maketitle

\noindent\emph{Source and licence.} Master equations with an individual noise on finite state graphs. Version arXiv:2605.06589v1 (CC BY 4.0), distributed under the Creative Commons Attribution 4.0 International licence, as recorded in the arXiv licence metadata for this exact version and shown on the abstract page https://arxiv.org/abs/2605.06589v1. Full rights notice: licenses/arxiv-2605.06589-CC-BY-4.0.txt.\par\smallskip

\noindent\textit{Editorial scope.} Counted unit: the Proposition whose source label is \texttt{prop:phi bd} in the article (source0.tex lines 655--659, source label \texttt{prop:phi bd}), delivered together with the article's own complete original proof of it (source0.tex lines 660--702, one \texttt{proof} environment of 43 lines). No mathematics is written, repaired or re-derived: statement and proof are the article's own text line for line. Required context retained uncounted: as:H lower bound (lines 596--601); as:g mu (lines 596--601); eq:mfg lambda (lines 618--623); lem:lasrylions (lines 626--630); lem:upper phi (lines 641--645). Every definition, display and label of those lines is retained unchanged. Omitted from this package and not redistributed: 1--595, 602--617, 624--625, 631--640, 646--654, 703--2070. Nothing inside a retained excerpt is omitted or altered: every retained body is delivered exactly as the article's lines are written, so no omission receipt of any kind accompanies this package. Citations: the retained excerpts carry no citation command at all, so no bibliography entry of the article is transcribed and no reference is redistributed. Reference adaptation: none was needed and none was made; every reference inside the delivered unit resolves inside it, so no unresolved reference or citation is produced. Printed slips kept literal: no printed word, symbol, hypothesis or number of the counted statement or of its proof was altered. Layout and class: the article's own class is \texttt{article} with packages including authblk, graphicx, amsfonts, amsopn, amsthm, amssymb, color, hyperref, soul, latexsym, amsgen, bigints; the wrapper is the lane's pinned \texttt{amsart} editorial wrapper at 10pt on A4, which restates the theorem-like environments the retained text needs and adds the article's own macro definitions that the retained text needs (\texttt{\textbackslash G}, \texttt{\textbackslash cP}), with extra wrapper packages (bm, bbm, dsfont, xspace, cancel, mathtools). The delivered numbering is the unit's own. Assets: no retained span contains a figure environment, a graphics-inclusion command, a PDF-page inclusion, a table or a TikZ picture; no image file is copied into this package, the package declares no asset dependency and no image file is redistributed with it. Licence: version arXiv:2605.06589v1 of this article is distributed under the Creative Commons Attribution 4.0 International licence, as recorded in the arXiv licence metadata for this exact version and shown on the abstract page https://arxiv.org/abs/2605.06589v1. A full rights notice is delivered at licenses/arxiv-2605.06589-CC-BY-4.0.txt. Characters: every retained excerpt is pure ASCII, so the delivered engine, XeLaTeX with its default Latin Modern text font, reports no missing character.
\begin{align}
 (B(\mu,p),p)&\geq (H(\mu,p),\mu)-C_1, \label{as:B coercivity}\\
 H_i(\mu,p)&\geq -C_1, \label{as:H lower bound}\\
 |B_{ij}(\mu,p)|&\leq (\mu^i+\mu^j)\,a(\mu^j/\mu^i,p), \label{eq:mobility-dominated flux}\\
 \max_{1\leq i\leq n,\,\rho\in[0,1]^n}|g_i(\rho)|&\leq C_1. \label{as:g mu}
\end{align}

\begin{equation}  \label{eq:mfg lambda}
    \begin{cases} \dot \phi(s)= H(\rho(s),\lambda\nabla_{\G}\phi(s))-\Delta_{\G}\phi(s)& s\in (t,T)\\
    \dot \rho(s)=\nabla_{\G}\cdot (B(\rho(s),\lambda\nabla_{\G}\phi(s)))+\Delta_{\G}\rho(s)& s\in (t,T)\\
    \phi(T)= g(\rho(T)), \quad \rho(t)=\mu.
    \end{cases}
\end{equation}

\begin{lemma}[Lasry--Lions estimate] \label{lem:lasrylions} Let $(\phi,\rho)$ be a classical solution to \eqref{eq:mfg lambda}. Then
\begin{equation}
\lambda ( \phi(T),\rho(T))\leq \lambda ( \phi(t),\mu)+2C_1(T-t).
\end{equation}
\end{lemma}

\begin{lemma}[Upper bound on $\phi$] \label{lem:upper phi} If $(\phi,\rho)$ is a classical solution to \eqref{eq:mfg lambda}, then
    \begin{equation}
        \max_{1\leq i \leq n} \phi_i(s) \leq  C_1(T-s)+\max_{1\leq i\leq n}  \phi_i(T), \quad s\in [t, T].
    \end{equation}
\end{lemma}

\begin{proposition}[Uniform bounds on $\lambda \phi$]\label{prop:phi bd} Let $\mu \in \cP_\epsilon(\G)$ and $t\in[0,T)$. For any classical solution $(\phi,\rho)$ to \eqref{eq:mfg lambda},
\begin{equation}
\epsilon \big\|\max_{1\leq i\leq n}| \lambda\phi_i|\big\|_{L^{\infty}(t,T)}\leq (4(T-t)+3)C_1.
\end{equation}
\end{proposition}

\begin{proof} By Lemma \ref{lem:upper phi}, we have
\begin{equation} \label{eq:lphi upper pf}
\max_{1\leq i\leq n} \lambda \phi_i(s) \leq \lambda (T-t+1)C_1, \quad s\in [t,T],
\end{equation}
so to prove the result it suffices to bound from below the function $s \mapsto \lambda\Phi(s)=(\lambda \phi(s),\textbf{1})=\sum_{i=1}^{n}\lambda\phi_i(s).$
%\[\lambda\Phi(s)=(\lambda \phi(s),\textbf{1})=\sum_{i=1}^{n}\lambda\phi_i(s).\]
Testing the first equation of \eqref{eq:mfg lambda} against \textbf{1}, using \eqref{as:H lower bound} and the fact that $(\Delta_{\G}\phi,\textbf{1})=0$, we have
$\dot \Phi(s)\geq\sum_{i}H_i \geq -C_1n.$
Therefore, noting that $1=\sum_{i}\mu^i\geq n\epsilon$, we obtain that
\begin{equation} \label{eq:dec16.2025.1}
\lambda \Phi(s)\geq \lambda \Phi(t)-\lambda C_1n(s-t)\geq \lambda \Phi(t)-\frac{C_1 }{\epsilon}(s-t),\quad s\in [t,T],
\end{equation}
so it suffices to bound $\lambda \Phi(t)$ from below. By Lemma \ref{lem:lasrylions}, we have
\begin{equation}
    -C_1 \leq (\lambda \phi(t),\mu)+2C_1(T-t) = (\lambda \phi(t)^+,\mu)+2C_1(T-t) -(\lambda\phi(t)^-,\mu).
\end{equation}
Moreover, by Lemma \ref{lem:upper phi} and \eqref{as:g mu},
\[
(\lambda \phi(t)^+,\mu)\le \max_{i}\lambda \phi_i(t)\le \max_i \phi_i(t)\le C_1(T-t)+\max_i\phi_i(T)\le C_1(T-t+1),
\]
and therefore
\begin{equation}
    -C_1 \leq (\lambda \phi(t),\mu)+2C_1(T-t) \leq C_1(2(T-t)+1) -(\lambda \phi(t)^-,\mu).
\end{equation}
Thus,
\begin{equation}\label{eq:dec16.2025.2}
    \lambda \Phi(t)^-\leq (\lambda \phi^-(t),\textbf{1})\leq \frac{1}{\epsilon}(\lambda \phi^-(t),\mu)\leq\frac{3C_1(T-t+1)}{\epsilon}.
\end{equation}
Recalling \eqref{eq:lphi upper pf}, we conclude that
\[
\lambda \phi_i(s)=  \lambda\Phi(s)-\sum_{k \not =i } \lambda \phi_k(s)
\geq \lambda\Phi(s) - (n-1)\max_k\lambda\phi_k(s),
\]
for all $i$. We use \eqref{eq:dec16.2025.1} and \eqref{eq:dec16.2025.2} to get that
\[
\lambda\Phi(s)\ge -\frac{3C_1(T-t+1)}{\epsilon}-\frac{C_1}{\epsilon}(s-t)\ge -\frac{C_1(4(T-t)+3)}{\epsilon}, \quad s\in[t,T],
\]
and hence, using again that $n\epsilon\le 1$,
\[
\lambda \phi_i(s)\ge -\frac{C_1(4(T-t)+3)}{\epsilon}-(n-1)C_1(T-t+1)\ge -\frac{C_1(5(T-t)+4)}{\epsilon},
\]
for all $i$ and $s \in [t,T].$ Together with \eqref{eq:lphi upper pf}, this proves the claim.
\end{proof}

\end{document}
